\chapter{Natural Effects}

Suppose we now observe a small effect of the anti-depressant on depression and are worried 
that a deleterious side effect associated with the anti-depressant, chronic headaches, might be interfering 
with the positive effects of the drug. In this situation, there are two possibilities: 

\begin{itemize}[noitemsep,topsep=0pt]
	\item The drug is effective at treating existing depressive symptoms, but the chronic pain from the side-effects 
	cancel out this benefit and provide a new mental health challenge. 
	\item The headaches are a minor mental health issue, but the drug is not effective at treating the original 
	symptoms. 
\end{itemize}

While the observed causal effect of the anti-depressant is the same in both situations, which of the two situations is 
operative has important implications for evaluating how promising the drug is. If we are in the first world,
a drug manufacturer might consider investing resources in reducing the side effects of an otherwise promising drug;
if we are in the second world, the anti-depressant does not seem like a promising treatment to invest in. 


\section{Effect Definitions}\label{subsection:natural_effect_defs}

We now define a mediator variable $M_{i}$, which represents a continuous measure of whether the intensity of chronic
headache pain a patient experienced during the study period. Let $\mathcal{M}$ represent the support of $M$. 
We also define potential mediators over $M_{i}$ and potential outcomes over $Y_{i}$:
\begin{itemize}[noitemsep,topsep=0pt]
	\item $M_{i}(0)$ --- The level of headache a patient would experience if they did not take anti-depressants.
	\item $M_{i}(1)$ --- The level of headache a patient would experience if they took anti-depressants.
	\item $Y_{i}(0,M_{i}(0))$ --- Depressive symptoms a patient would experience if they were untreated and had 
	the level of headache they would experience if they were not treated. 
	\item $Y_{i}(1,M_{i}(1))$ --- Depressive symptoms a patient would experience if they were treated and had 
	the level of headache they would experience if they were treated. 
	\item $Y_{i}(1,M_{i}(0))$ --- Depressive symptoms a patient would experience if they were treated but had 
	the level of headache they would experience if they were untreated.
	\item $Y_{i}(0,M_{i}(1))$ --- Depressive symptoms a patient would experience if they were untreated but had 
	the level of headache they would experience if they were treated.
\end{itemize}

If a unit does not receive treatment, we observe their realized mediators $M_{i}(0)$ and outcomes $Y_{i}(0,M_{i}(0))$;
if they are treated, we observe $M_{i}(1)$ and $Y_{i}(0,M_{i}(1))$. Note that we can never observe $Y_{i}(1,M_{i}(0))$ or
$Y_{i}(0,M_{i}(1))$ because a unit can only be assigned a treatment level once, but evaluating $Y_{i}(t',M_{i}(t))$ for 
$t' \neq t$ would require knowing both the unit's outcome under treatment level $t'$ and their mediator value under
treatment level $t$~\citep{robinsgreenland1992}. There does not exist an experimental procedure to identify 
these cross-world counterfactual~\citep{pearl2001} without making an additional set of hard-to-satisfy 
assumptions~\citep{imaietal2013}. Having defined this expanded set of potential mediators and outcomes, we
now define a set of estimands that can help define the path-specific causal effects we are interested in. Note that 
these effects are also defined at the unit-specific level but are never identified at the level, so the 
definitions we provide correspond to the corresponding population average estimands for these effects. Note that
we define average effects across the entire study population as opposed to a treated or untreated subset of the 
population. 

\begin{itemize}[noitemsep,topsep=0pt]
	\item $\mathbb{E}[Y_{i}(1,M_{i}(0)) - Y_{i}(0,M_{i}(0))]$ --- The \textit{pure direct effect} (PDE)
	is the average causal effect that receiving treatment has on the outcome, keeping the level of the mediator constant 
	for each unit at the level of the mediator that would have been realized had the unit not received treatment. 
	\item $\mathbb{E}[Y_{i}(1,M_{i}(0)) - Y_{i}(0,M_{i}(0))]$ --- The \textit{total direct effect} (TDE)
	is the average causal effect that receiving treatment has on the outcome, keeping the level of the mediator constant 
	for each unit at the level of the mediator that would have been realized had the unit received treatment. 
	\item $\mathbb{E}[Y_{i}(0,M_{i}(1)) - Y_{i}(0,M_{i}(0))]$ --- The \textit{pure indirect effect} (PIE)
	is the average causal effect that changes in the mediator that are induced from changes in the treatment level have
	on the outcome, evaluated at the untreated level of the treatment.  
	\item $\mathbb{E}[Y_{i}(1,M_{i}(1)) - Y_{i}(1,M_{i}(0))]$ --- The \textit{total indirect effect} (TIE)
	is the average causal effect that changes in the mediator that are induced from changes in the treatment level have
	on the outcome, evaluated at the treated level of the treatment.  
	\item $\mathbb{E}[Y_{i}(1,m^{*}) - Y_{i}(0,m^{*})]$ --- The \textit{controlled direct effect} (CDE) at $M=m^{*}$ is the 
	average causal effect that receiving treatment has on the outcome, keeping the level of the mediator fixed at a 
	particular level $m$ for all units. There is a (potentially different) CDE defined at all possible values of the 
	mediator. 
	\item $\mathbb{E}[Y_{i}(1,M_{i}(1)) - Y_{i}(0,M_{i}(0))]$ --- The \textit{total effect} (TE) is 
	the average causal effect of a change in the treatment on the outcome. This is equal to the ATE and is identified
	under the same assumptions. 
\end{itemize}

The pure and total direct and indirect effects are known in the literature as 
\textit{natural effects}~\citep{robinsgreenland1992,pearl2001} because 
they attempt to measure causal effects while allowing for the "natural relationship"~\citep{pearl2001} between the treatment 
and the mediator to be expressed.

In the running example, the \textit{natural direct effect} measures what the average effect of 
the anti-depressant on depression would be in a counterfactual world in which the drug had no side effect of headache. 
This would mean that each patient would experience the same patient-specific, "natural" level of headache 
that they would have if they had not taken the drug. Thus, this effect
gives the effect of the anti-depressant on depression without any of the effects that result from the treatment's 
side effects. The PDE operationalizes this idea by comparing a world in which the patient does not receive treatment with 
another where the patient receives treatment but it is imagined that the mediator is held at the level it would have
attained in the first world. 

The crucial difference between the pure and total direct effects and the controlled direct effect is that the level of the 
mediator is allowed to vary between units and is set to an unobservable "natural" level in the definition of the pure and
total effects while it is externally fixed 
for all units at some level $M = m^{*}$ for all units for the controlled effect. These two effects are equivalent when there 
is no interaction between the treatment and mediator levels that affects the outcome~\citep{pearl2001}, and the controlled 
effect will be the same at all levels of $M$ in this case. As described in~\citep{pearl2001}, the 
controlled effect is useful for prescriptive goals, such as a question of how to make a treatment regime most effective 
when the mediator can also be manipulated; natural effects are descriptive and give insight into which causal mechanisms 
are operative in the observation of a total causal effect. There is no equivalent controlled 
indirect effect~\citep{pearl2001}. 

The natural indirect effect measures the average effect of the change in chronic headache severity 
induced by taking the anti-depressant on depression. In this situation, it is imagined that the treatment administered
to the patient has no effect on depression other than through its effect on headache. The idea is operationalized by
a comparison between a world in which the patient either does (total indirect effect) or does not (pure indirect effect)
receive the treatment and a world in which they still receive the same treatment level as in the first world, 
but it is imagined that their mediator level is taken from a world in which they received the other treatment level. 


\section{Pure Versus Total Natural Effects}\label{subsection:pure_total_natural_effects}

Following~\cite{pearl2001}, many papers in the literature do not make a distinction between a "pure" and "total"
effect, instead defining a natural direct effect in an equivalent way to how we have defined the 
pure direct effect and a natural indirect effect in an equivalent way to how we have defined the total indirect effect.
Defining direct and indirect effects in this way allows one to decompose the total effect into these two effects, which
is naturally of interest given that the development of natural effects was motivated in part by an attempt to attribute 
responsibility for a causal effect to different causal mechanisms~\citep{pearl2001}.~\cite{pearl2001} gives an exact 
decomposition of the total effect in terms of (in our terms) the pure direct effect and the total indirect effect:

\begin{equation}
\begin{aligned}
	\mathbb{E}[Y_{i}(1,M_{i}(1)] - \mathbb{E}[Y_{i}(0,M_{i}(0))] = 
	\{\mathbb{E}[Y_{i}(1,M_{i}(0))] - \mathbb{E}[Y_{i}(0,M_{i}(0))]\}\\ 
	+ \{\mathbb{E}[Y_{i}(1,M_{i}(1))] - \mathbb{E}[Y_{i}(1,M_{i}(0))]\}
\end{aligned}
\end{equation}

It is easy to see that the same trick of adding and subtracting a term from the total effect gives an alternative
exact decomposition in terms of the total direct effect and the pure indirect effect. However, the decomposition
in terms of the pure direct and total indirect effects is the more common one in the literature~\citep{vanderweele2014}. 
As detailed in
~\cite{robinsgreenland1992}, the pure and total natural effects are not in general equivalent. In these situations, 
the choice of which decomposition "correctly" represents the direct and indirect effects is unclear.

To illustrate this point,~\cite{robinsgreenland1992} analyzes an example with a binary mediator and binary
outcome from the perspective of the sufficient-component causes framework~\citep{rothman1976}. We briefly assume 
that the mediator is a binary indicator for the presence of chronic headache, while
the outcome is the presence of a dichotomized depression indicator.~\cite{robinsgreenland1992} develops a typology
of patient types based on hypothetical answers for that patient to the following questions:

\begin{enumerate}[noitemsep,topsep=0pt]
\item Headache would occur if the drug is taken.
\item Headache would occur if the drug is not taken.
\item Depression would occur if the drug is taken but headache is not experienced.
\item Depression would occur if the drug is taken and headache is experienced.
\item Depression would occur if the drug is not taken but headache is not experienced.
\item Depression would occur if the drug is not taken and headache is experienced.
\end{enumerate}

We also assume for now that patients who would develop headache when not taking the anti-depressant but do not in the 
drug's absence exist. Given this typology, we can say that for patients for whom 5 and 6 are true and 3 and 4 are
false, we can say that the drug has had a direct effect on the patient's depression: depression occurs for the 
patient only in the absence of the drug, regardless of whether they experience headache. 
For patients for whom 1, 4, and 6 are true and 2, 3, and 5 are false (here, the additional condition that 1 is 
true and 2 is false ensure that the treatment actually has an effect on the level of headahce), we can say
that the drug has had an indirect effect on depression: headache occurs only in the presence of the drug,
and depression occurs only in the presence of headache, regardless of whether the drug is taken. 

Now consider a patient for whom headache occurs only if the anti-depressant is taken. Additionally, they 
experience depression only if the drug is taken \textit{and} they experience headache, but do not experience
depression otherwise. For this patient, we might say that the interaction between the anti-depressant and 
headache leads to depression. However, it is unclear whether such an effect should be attributed to a direct or
an indirect effect: a necessary condition for the depression seems to be both the presence of the drug and 
the effect that the drug has on headache. In the definitions we have laid out, such a patient's status would be
attributed to the total direct or indirect effect but not the pure direct or indirect effect. This is why the
total causal effect decomposes to a sum of one total and one pure effect: summing two pure effects would not
count the effect on patients for whom the treatment and mediator statuses interact, while summing total effects 
would double count these interaction effects. 

~\cite{vanderweele2013} approaches the ambiguity in the definition of pure and total natural effects by proposing
a 3-way decomposition of the total effect into the pure direct effect, the pure indirect effect, and a third
term that can be interpreted as the causal effect of interactions between the treatment and mediator:

\begin{equation}
\begin{aligned}
	\mathbb{E}[Y_{i}(1,M_{i}(1))] - \mathbb{E}[Y_{i}(0,M_{i}(0))] = 
	\{\mathbb{E}[Y_{i}(1,M_{i}(0))] - \mathbb{E}[Y_{i}(0,M_{i}(0))]\}\\ 
	+ \{\mathbb{E}[Y_{i}(1,M_{i}(1))] - \mathbb{E}[Y_{i}(1,M_{i}(0))]\}\\
	= \{\mathbb{E}[Y_{i}(1,M_{i}(0))] - \mathbb{E}[Y_{i}(0,M_{i}(0))]\}
	+ \{\mathbb{E}[Y_{i}(0,M_{i}(1))] - \mathbb{E}[Y_{i}(0,M_{i}(0))]\} + \\
	\{(\mathbb{E}[Y_{i}(1,M_{i}(1))] - \mathbb{E}[Y_{i}(1,M_{i}(0))]) - 
	(\mathbb{E}[Y_{i}(0,M_{i}(1))] - \mathbb{E}[Y_{i}(0,M_{i}(0))])\}
\end{aligned}
\end{equation}

% Note that "pure" vs. "total" is also unclear, since "pure" does not necessarily mean the absence of interaction
% the interaction term captures only the difference between the pure and total effects; might be different at 
% different values of treatment; an example where treatment is not a drug with dosage but instead something like two 
% qualitatively different treatment regimes (e.g. therapy vs. anti-depressant) could drive this point home

~\cite{vanderweele2014} extends this result to a 4-way decomposition that decomposes the pure direct effect into 
the sum of the controlled direct effect and the component of the total effect that is due to interaction but not 
mediation (in our example, this would correspond to patients who would experience headache no matter what, but 
develop depression only when both the anti-depressant and headache are present). We do not focus further on this 
more granular decomposition in this paper. However, an analogue to the 3-way decomposition for interventional
effects exists, and this decomposition provides the conceptual foundation for our discussion of interventional 
effects with multiple causally ordered mediators in Section~\ref{chapter:sample}. 

\section{Identification Assumptions}

In order to identify most of the effects discussed in section~\ref{subsection:natural_effect_defs}, a number of 
additional identification assumptions are required. The one exception is the total effect, which is both mathematically 
and conceptually equivalent to the average treatment effect in the "simple" causal inference problem with no mediators 
and is thus identified under only the SUTVA and conditional ignorability assumptions 
from Section~\ref{subsection:natural_effect_defs}. 

% -------
% Think about how to present these assumptions; they are only mostly relevant for the controlled effect, and should we
% be presenting the controlled effect here at all? Or can we actually move this definition to the next section? The 
% argument for keeping it here is the strong discussion of Pearl's descriptive vs. prescriptive distinction, which 
% requires the introduction of the controlled effect here. So maybe we can still introduce these assumptions here. 
We now introduce a generalized SUTVA assumption that gives consistency of both the treatment and mediator and 
no interference between units~\citep{imaietal2013}. This background technical assumption is assumed in most of the 
causal mediation analysis literature~\citep{pearl2001,imaietal2010a}.

\begin{assumption}[Consistency]\label{assm:composition}
$M_{i}(a) = M_{i}$ if unit $i$ receives treatment level $a$. 
$Y_{i}(a,m) = Y_{i}$ if unit $i$ receives treatment level $a \in \mathcal{A}$ and 
has mediator level
\end{assumption}

\begin{assumption}[Composition]\label{assm:composition}
$M_{i}(a) = m$. $Y_{i}(a) = Y_{i}(a,m)$ if $M_{i}(a) = m$. 
\end{assumption}

The consistency portion of the assumption, which is the more contentious assumption in this context~\citep{imaietal2013}, 
states that each combination of treatment and mediator for a unit maps to one potential outcome for that unit, 
regardless of how those values are set. In particular, it is assumed that an intervention to set the mediator to the 
natural level of the mediator given a treatment $a$ will lead to the same potential outcome for a unit as if that 
mediator had not been set but instead allowed to attain its natural level. 

% --------

We also introduce a series of ignorability assumptions from the literature~\cite{pearl2001,imaietal2010a} that generalize
Assumption~\ref{assm:ignorability}.

\begin{assumption}[Sequential Ignorability]\label{assm:seq_ignorability} 
\indent
\begin{enumerate}[noitemsep,topsep=0pt,label=(\roman*)]
	\item $Y_{i}(a,m) \indep A_{i} | C_{i}$ for all $a \in \mathcal{A}$, $m \in \mathcal{M}$. 
	\item $M_{i}(a) \indep A_{i} | C_{i}$ for all $a \in \mathcal{A}$. 
	\item $Y_{i}(a,m) \indep M_{i} | C_{i}$ for all $a \in \mathcal{A}$, $m \in \mathcal{M}$. 
	\item $Y_{i}(a,m) \indep M_{i}(a^{*}) | A_{i} = a, C_{i} = c$ for all $a,a^{*} \in \mathcal{A}$, $m \in \mathcal{M}$. 
\end{enumerate}
where $0 < \mathbb{P}[A_{i} = a | C_{i}] < 1$ and $0 < \mathbb{P}[M_{i} = m | A_{i} = a, C_{i}]$
 for all $a \in \mathcal{A}$, $m \in \mathcal{M}$. 
\end{assumption}

Parts (i)-(ii) of Assumption~\ref{assm:seq_ignorability} state that conditional on some covariates $C_{i}$, the 
potential mediators $M_{i}(a)$ and outcomes $Y_{i}(a,m)$ are independent of treatment assignment $a$. Such an 
assumption is guaranteed to hold in randomized experiments where the treatment is randomly assigned. Additionally, 
(iii) is guaranteed to hold in randomized experiments where the mediator is randomly assigned in addition to the 
treatment. In the running example, an experiment that both randomly prescribed the anti-depressant and a dose of a 
(hypothetical) headache medication that sets a level of headache but has no other effects would fulfill the 
first part of the sequential ignorability assumption. In observational studies, similar strategies to those pursued
in observational studies in simple causal inference problems --- for example, collecting information on a number
of relevant confounding variables --- can be pursued to fulfill this assumption, except that confounders of both 
the treatment-outcome and the treatment-mediator relationship need to be controlled for~\citep{imaietal2010b}.  

Part (iv) of Assumption~\ref{assm:seq_ignorability}, commonly referred to the cross-world independence assumption, 
requires that the potential mediators and the potential outcomes are independent of each other. 
An alternative interpretation is that, among units that 
share a set of covariates $C_{i}$ and treatment assignment $a$, the natural level of the mediator for each unit 
can be treated as if it were randomized~\citep{imaietal2010a}. Unlike the first ignorability assumption, there does not 
exist an experiment that can fulfill this assumption without difficult-to-satisfy 
auxilliary assumptions~\citep{imaietal2013}. A representative violation of this condition would occur if there were a 
post-treatment confounder $L$ of the relationship between $M$ and $Y$.

The hypothetical experiment outlined above would not fulfill the cross-world independence assumption 
because the randomization of the mediator does not give any information about the potential mediators
given treatment. In general, the cross-world independence assumption involves an assumption about the 
independence of $Y_{i}(a,m)$ and $M_{i}(a^{*})$, which can never be seen for the same unit because it would involve
seeing the mediator at one treatment level and the outcome at another. Thus, it is considered to be empirically 
unfalsifiable~\citep{vanderweelevansteelandt2009, vanderweeleetal2014}. 

% -------
% More discussion of controlled direct effect
The controlled direct effect is identifiable under generalized SUTVA and the first part of the sequential
ignorability assumption~\citep{hernan2024}. 
% -------

All forms of the natural direct and indirect effects require 
that the cross-world independence assumption holds in addition to (non-generalized) SUTVA and the first part of the 
sequential ignorability assumption~\citep{pearl2001, imaietal2010a}. The implied DAG representing
the relationships between $C$, $A$, $M$, and $Y$ under these identification assumptions is given in 
Figure~\ref{fig:one_mediator_dag}. 


\begin{figure}[ht]
\centering
\begin{tikzpicture}

  % Define nodes
  \node[obs]                               (A) {$\mathbf{A}$};
  \node[obs, right=of A, yshift=1.5cm] (M) {$\mathbf{M}$};
  \node[obs, right=2.5cm of A]  (Y) {$\mathbf{Y}$};
  \node[obs, left=of A]  (C) {$\mathbf{C}$};

  % Connect the nodes
  \edge {A,M} {Y} ; %
  \draw[->] (C) to[out=315, in=225] (Y); %
  \edge {C,A} {M} ; %
  \edge {C} {A} ; %

\end{tikzpicture}
\caption[]{Mediation analysis with one mediator}
\label{fig:one_mediator_dag}
\end{figure}

The sequential ignorability assumption is slightly stronger than is necessary for the identification of 
natural effects. For example, in situations where there is no interaction between the treatment and the mediator,
the controlled direct effect is equal to the natural direct effect, and thus natural effects are identified without
the cross-world independence assumption~\citep{robins2003}. Additionally, it is possible to slightly weaken this no
interaction assumption and still identify natural effects~\citep{petersenetal2006}. However, these alternative conditions
for identification are not generally regarded as being meaningfully weaker than sequential ignorability in practical 
situations are more difficult to interpret~\citep{petersenetal2006, imaietal2010a}. 

\section{Multiple Mediators}

In addition to being empirically unfalsifiable, a key limitation of the cross-world independence assumption is that it 
concretely rules out the existence of treatment-induced confounders. 
These are variables that might be affected by the treatment and themselves confound the relationship between the mediator 
of interest and the outcome. The natural direct and indirect effects can in general not be identified in these 
situations without alternative strong assumptions, even if these treatment-induced confounders can be 
measured~\citep{vanderweelevansteelandt2009}. Figure~\ref{fig:exposure_induced_confounder} shows an example of such a 
situation.

\begin{figure}
\centering
\begin{tikzpicture}

  % Define nodes
  \node[obs] (T) {$\mathbf{T}$};
  \node[obs, right=of T, yshift=1.5cm] (L) {$\mathbf{L}$};
  \node[obs, right=2cm of T] (M) {$\mathbf{M}$};
  \node[obs, right=1cm of M]  (Y) {$\mathbf{Y}$};
  \node[obs, left=of T]  (C) {$\mathbf{C}$};

  % Connect the nodes
  \edge {T,L} {M} ; %
  \edge {M} {Y} ; %
  \draw[->] (T) to[out=-30, in=-150] (Y); %
  \draw[->] (L) to[out=0, in=120] (Y); %
  \draw[->] (C) to[out=300, in=240] (Y); %
  \draw[->] (C) to[out=315, in=225] (M); %
  \edge {C,T} {L} ; %
  \edge {C} {T} ; %

\end{tikzpicture}
\caption[]{One mediator, exposure-induced confounding}
\label{fig:exposure_induced_confounder}
\end{figure}

The assumption of no treatment-induced confounding is difficult to rule out and limits the space of situations
were natural mediational effects can be credibly identified. Crucially for the focus of this dissertation, it also means
that in situations with multiple mediators that are not assumed to be independent of each other,
traditional direct and indirect effects cannot be identified~\citep{avinetal2005,imaiyamamoto2013}. To see this, consider 
Figure~\ref{fig:exposure_induced_confounder} with $L$ as a mediator of interest as opposed to a 
confounder. 

In causal structures like the one represented in 
Figure~\ref{fig:exposure_induced_confounder},~\cite{avinetal2005} showed that it is possible to 
decompose the total effect into three path-specific natural effects, $A \rightarrow Y$, 
$A \rightarrow M \rightarrow Y$, and the sum of $A \rightarrow L \rightarrow M \rightarrow Y$ and
$A \rightarrow L \rightarrow Y$ given an expanded set of sequential ignorability assumptions is fulfilled. 
Relatedly, one can treat the mediators as a single mediator and decompose the total effect into effects through
$A \rightarrow Y$ versus all other pathways~\citep{vanderweeleetal2014}. However, neither option retains the notion
of natural direct and indirect effects that we are interested in. In summary, if we are interested in problems with 
multiple mediators that exist in some known or unknown causal ordering with each other, the goal of using natural 
effects to attribute portions of an observed total causal effect to effects through specific causal pathways will 
generally not be a fruitful one. 

% Notes:
% - Add a discussion of the fact that an exact decomposition of the total effect is of interest given the 
% descriptive/attributional goals of natural effects; less important in the prescriptive controlled/interventional effects case